\documentclass[10pt,a4paper]{amsart}
\usepackage[margin=19mm]{geometry}
\usepackage{amsmath,amssymb,mathtools}
\usepackage[scr=rsfs,cal=euler]{mathalfa}
\usepackage{xfrac}
\usepackage[unicode,hidelinks,bookmarks=false]{hyperref}
\newcommand{\ie}{i.e.\ }
\newcommand{\cf}{cf.\ }
\newcommand{\resp}{respectively\ }
\newcommand{\Subsection}{\subsection}
\providecommand{\nobreakdash}{\nobreak\hbox{-}}
\setlength{\emergencystretch}{2em}
\multlinegap0pt
\usepackage{amssymb}
\newtheorem{theorem}{Theorem}
\newtheorem{lemma}{Lemma}
\newtheorem{proposition}{Proposition}
\usepackage{cleveref}
\crefname{theorem}{Theorem}{Theorems}
\crefname{lemma}{Lemma}{Lemmas}
\crefname{proposition}{Proposition}{Propositions}
\newcommand{\Arch}{P_{\mathrm{Arch}}}
\newcommand{\Ical}{\mathcal I}
\newcommand{\e}{\mathrm e}
\newcommand{\mb}[1]{\mathbf{#1}}
\newcommand{\vol}{\operatorname{vol}}
\title{Local quadratic isoperimetric stability of the truncated octahedron among parallelohedra}
\author{Lark Song}
\date{Source: arXiv:2609.03093v1 (CC BY 4.0)}
\begin{document}
\maketitle

\noindent\emph{Source and licence.} Local quadratic isoperimetric stability of the truncated octahedron among parallelohedra. Version arXiv:2609.03093v1 (CC BY 4.0), distributed by arXiv under the Creative Commons Attribution 4.0 International licence, http://creativecommons.org/licenses/by/4.0/, as recorded in the arXiv licence metadata for this exact version and shown on the abstract page https://arxiv.org/abs/2609.03093v1. Full licence notice: \texttt{licenses/arxiv-2609.03093-CC-BY-4.0.txt}.\par\smallskip

\noindent\textit{Editorial scope.} Counted unit: the article's theorem thm:main (source0.tex lines 111--118). The article's complete original proof of it is delivered with it (source0.tex lines 635--730, one \texttt{proof} environment of 96 lines, which the article places away from the statement). No mathematics is written, repaired or re-derived: the statement and the proof are the article's own lines, in the article's own notation, and no retained line is substituted or re-flowed.  Retained uncounted as the source-local context the proof needs: lines 101--103; lines 329--332; lines 380--386; lines 495--497; lines 504--510.  Nothing was omitted from any retained span, so the package carries no omission record.  No retained line contains a citation, so the wrapper carries no bibliography and no bibliography entry.  No retained line contains a figure environment, an image-inclusion command or drawing code, so no image is delivered and the package declares no asset dependency.  Editorial formatting only, in addition: an AMS \texttt{amsart} preamble that restates the article's own declarations and macros used by the retained lines, namely the environments \texttt{theorem}, \texttt{lemma}, \texttt{proposition} on separate counters, together with the standard packages those lines need.  The retained environments print in this document as Theorem 1, Lemma 1, Proposition 1 in order of appearance, whereas the article prints the same items with the numbering of its own full document.  The complete native source of the article as distributed in the arXiv e-print for this exact version is bundled unchanged as \texttt{sources/209264/source0.tex}.
\begin{equation}\label{eq:distance}
 d_{\mathrm{sim}}(P,\Arch):= \inf_{Q\in\mathrm{SO}(3)}d_H(Q\widehat P,\Arch).
\end{equation}

\begin{lemma}[openness of the truncated-octahedral class]\label{lem:open-class}
There is a Hausdorff neighborhood of $\Arch$ in $\mathcal P_3$ that consists entirely of
truncated-octahedral parallelohedra.
\end{lemma}

\begin{proposition}[Hausdorff-compatible quotient chart]\label{prop:topology}
The map $h\mapsto[\Phi(h)]$ induces a local homeomorphism from a neighborhood of $[0]$ in
$(\mathcal W_0\oplus\mathcal G_0)/S_4$ onto a neighborhood of $[\Arch]$ in the
similarity quotient of the truncated-octahedral class.  Equivalently, $[P_k]\to[\Arch]$
if and only if, after relabeling the six zones, the product-normalized representatives are
$\Phi(h_k)$ with $h_k\to0$.
\end{proposition}

\begin{equation}\label{eq:critical}
 D\Psi(0)=0.
\end{equation}

\begin{equation}\label{eq:decomposition}
 \mathcal W_0\simeq E\oplus T,
 \qquad
 \mathcal G_0\simeq E\oplus T,
 \qquad
 \mathcal W_0\oplus\mathcal G_0\simeq2E\oplus2T.
\end{equation}

\begin{theorem}[quadratic local stability]\label{thm:main}
There exist a Hausdorff neighborhood $\mathcal U$ of $\Arch$ in $\mathcal P_3$ and a constant $c_{\mathrm{sim}}>0$ such that
\begin{equation}\label{eq:intrinsic-stability}
 \Ical(P)-\Ical(\Arch)\geq
 c_{\mathrm{sim}}d_{\mathrm{sim}}(P,\Arch)^2.
\end{equation}
For $P\in\mathcal U$, equality holds if and only if $P$ is similar to $\Arch$. 
\end{theorem}

\begin{proof}[Completion of the proof of \Cref{thm:main}]
Their first diagonal entries are positive, and
\[
 \det H_E=\frac{9(2\sqrt3-3)}{2^{23/3}}>0,\qquad
 \det H_T=\frac{5+2\sqrt3}{12\cdot 2^{8/3}}>0.
\]
Hence $H_E$ and $H_T$ are positive definite.  By \eqref{eq:decomposition},
$D^2\Psi(0)$ is positive definite on all ten tangent directions.  Thus the
origin is a nondegenerate strict minimum for $\Psi$ on the ten-dimensional labeled
similarity slice.
Analyticity, \eqref{eq:critical} and Taylor's theorem give
\[
 \Psi(h)=\Psi(0)+\frac12D^2\Psi(0)[h,h]+o(|h|^2).
\]
After shrinking the chart, there are $c,\varepsilon>0$ such that
\begin{equation}\label{eq:chart-gap}
 \Psi(h)-\Psi(0)\geq c|h|^2
 \qquad(|h|<\varepsilon).
\end{equation}
The deficit therefore vanishes only at $h=0$ in the labeled chart.  Restoring scale, rotation, and finite relabeling yields precisely the similarity class of $\Arch$.
The support function is
\begin{equation}\label{eq:support}
 h_{\Phi(\ell,\sigma)}(u)=\frac12\sum_{i<j}\e^{\ell_{ij}}
 \left|u\cdot \e^{-\sigma/2}\mb n_{ij}^{\times}\right|,
 \qquad u\in S^2.
\end{equation}
For volume normalization, let
$t(h):=(V_0/V(\e^\ell))^{1/3}$.  Since $\vol(\Arch)=V_0$, the normalization in
\eqref{eq:distance} reads $$\widehat{\Phi(h)}=t(h)\Phi(h).$$
The factor $t$ is analytic and $t(0)=1$.  Fix $r>0$ small enough that
the closed ball $\{|h|\leq r\}$ is contained in the coordinate chart.
Since the scalar exponential and the matrix exponential
are $C^1$,
there is $L_r>0$ such that, whenever $|h|\leq r$,
\[
 |\e^{\ell_{ij}}-1|\leq L_r|\ell_{ij}|,\qquad
 \|\e^{-\sigma/2}-I\|_F\leq L_r\|\sigma\|_F,\qquad
 |t(h)-1|\leq L_r|h|.
\]
Moreover, $\e^{\ell_{ij}}$ and $\|\e^{-\sigma/2}\|_F$ are uniformly
bounded on this ball.  Since $\Phi(0)=\Arch$, the inequality
$\bigl||a|-|b|\bigr|\leq|a-b|$ gives, uniformly for $u\in S^2$,
\[
\begin{aligned}
 |h_{\Phi(h)}(u)-h_{\Arch}(u)|
 &\leq \frac12\sum_{i<j}
 \Bigl(
  |\e^{\ell_{ij}}-1|
  \left|u\cdot\e^{-\sigma/2}\mb n_{ij}^{\times}\right|
  +
  \left|u\cdot
  (\e^{-\sigma/2}-I)\mb n_{ij}^{\times}\right|
 \Bigr)                                                     \\
 &\leq C_1\left(\sum_{i<j}|\ell_{ij}|+\|\sigma\|_F\right)
 \leq C_2|h|
\end{aligned}
\]
for constants $C_1,C_2>0$.  The same bounds show that
$\|h_{\Phi(h)}\|_\infty$ is uniformly bounded for $|h|\leq r$.
Consequently,
\[
\begin{aligned}
 \|h_{\widehat{\Phi(h)}}-h_{\Arch}\|_\infty
 &=
 \|t(h)h_{\Phi(h)}-h_{\Arch}\|_\infty                         \\
 &\leq
 |t(h)-1|\|h_{\Phi(h)}\|_\infty
 +\|h_{\Phi(h)}-h_{\Arch}\|_\infty              \leq C_0|h|
\end{aligned}
\]
after increasing the constant if necessary.  Hausdorff distance is the
uniform distance between support functions.  Hence, choosing the identity
rotation in the defining infimum,
\begin{equation}\label{eq:support-bound}
 d_{\mathrm{sim}}(\Phi(h),\Arch)
 \leq d_H(\widehat{\Phi(h)},\Arch)
 =\|h_{\widehat{\Phi(h)}}-h_{\Arch}\|_\infty
 \leq C_0|h|.
\end{equation}  
Take $c_{\mathrm{sim}}:=c/(2C_0^2)$. 
This choice makes the asserted inequality strict away from the Archimedean similarity class.  
By \Cref{lem:open-class} and \Cref{prop:topology}, after shrinking the Hausdorff neighborhood
$\mathcal U$, every $P\in\mathcal U$ is truncated-octahedral and its similarity class has
a lift $h$ with $|h|<\varepsilon$.  Any two such lifts differ by the orthogonal $S_4$
action, so $|h|$ is independent of the lift.  Similarity
invariance, \eqref{eq:chart-gap}, and \eqref{eq:support-bound} give
\[
 \Ical(P)-\Ical(\Arch)=\Psi(h)-\Psi(0)
 \geq c|h|^2
 \geq \frac{c}{C_0^2}d_{\mathrm{sim}}(P,\Arch)^2
 =2c_{\mathrm{sim}}d_{\mathrm{sim}}(P,\Arch)^2.
\]
This proves \eqref{eq:intrinsic-stability}.  If $P$ is not similar to $\Arch$, then
$d_{\mathrm{sim}}(P,\Arch)>0$, and the last bound is strictly larger than the right-hand
side of \eqref{eq:intrinsic-stability}.  If $P$ is similar to $\Arch$, both sides vanish.
\end{proof}

\end{document}
